% !TEX root = 第六章_相似标准型(答案版).tex
\setcounter{chapter}{5}
\pagestyle{fancy}

\chapter{相似标准型}

本章是上册``线性变换''之后的进一步深化,核心目标是把线性变换按照
最小多项式,准素分解,循环分解与三种因子逐步化为有理标准形和
Jordan 标准形,并在此基础上讨论矩阵幂,可对角化,幂零矩阵,方阵函数
以及矩阵可交换问题.

\section{准素分解与循环分解(Jordan 标准型 1)}

\subsection{零化多项式,最小多项式与特征多项式}

\begin{definition}
设 $A$ 为方阵,若 $g(\lambda)\in F[\lambda]$ 且
\[
g(A)=O,
\]
则称 $g(\lambda)$ 为 $A$ 的\strongcn{零化多项式}.
\end{definition}

\begin{definition}
$A$ 的次数最低的首一零化多项式称为 $A$ 的\strongcn{最小多项式},
记为 $m_A(\lambda)$.
\end{definition}

\begin{theorem}[\markstar]
最小多项式存在且唯一,并且整除 $A$ 的任意零化多项式.
\end{theorem}

\begin{theorem}[\markstar\ Cayley--Hamilton 定理]
若
\[
f_A(\lambda)=\det(\lambda I-A)
\]
是 $A$ 的特征多项式,则
\[
f_A(A)=O.
\]
\end{theorem}

\begin{theorem}[\markstar]
在复数域上,$A$ 的特征多项式与最小多项式具有相同的根集,
但各根的重数可以不同.
\end{theorem}


%例题6.1.1
\begin{example}
求下列方阵的最小多项式:
\begin{enumerate}
    \item
    \[
    A_1=
    \begin{pmatrix}
    5&-6&-6\\
    -1&4&2\\
    3&-6&-4
    \end{pmatrix};
    \]
    \item
    \[
    A_2=
    \begin{pmatrix}
    3&1&-1\\
    2&2&-1\\
    2&2&0
    \end{pmatrix};
    \]
    \item
    \[
    A_3=
    \begin{pmatrix}
    0&-1\\
    1&0
    \end{pmatrix};
    \]
    \item Jordan 块 $J_n(\lambda_0)$.
\end{enumerate}
\end{example}
\exampleblank

\begin{solution}
前两个矩阵的特征多项式均为
\[
(\lambda-1)(\lambda-2)^2.
\]
对 $A_1$ 直接验证
\[
(A_1-I)(A_1-2I)=O,
\]
故
\[
\boxed{m_{A_1}(\lambda)=(\lambda-1)(\lambda-2)}.
\]
而 $A_2$ 的特征值 $2$ 的几何重数为 $1$,故最小多项式中
$(\lambda-2)$ 的指数必须为 $2$:
\[
\boxed{m_{A_2}(\lambda)=(\lambda-1)(\lambda-2)^2}.
\]

对第三问,
\[
A_3^2=-I,
\]
故
\[
\boxed{m_{A_3}(\lambda)=\lambda^2+1}.
\]

最后,
\[
(J_n(\lambda_0)-\lambda_0I)^n=O,
\qquad
(J_n(\lambda_0)-\lambda_0I)^{n-1}\neq O,
\]
所以
\[
\boxed{m_{J_n(\lambda_0)}(\lambda)=(\lambda-\lambda_0)^n}.
\]
\end{solution}


%例题6.1.2
\begin{example}[\marktwostar]
设
\[
f(\lambda)=\lambda^n+c_{n-1}\lambda^{n-1}+\cdots+c_1\lambda+c_0,
\]
其友矩阵为
\[
C(f)=
\begin{pmatrix}
0&0&\cdots&0&-c_0\\
1&0&\cdots&0&-c_1\\
0&1&\cdots&0&-c_2\\
\vdots&\vdots&\ddots&\vdots&\vdots\\
0&0&\cdots&1&-c_{n-1}
\end{pmatrix}.
\]
证明 $C(f)$ 的最小多项式和特征多项式均为 $f$.
\end{example}
\exampleblank

\begin{solution}
令 $e_1,\ldots,e_n$ 为标准基,则
\[
e_1,\ C(f)e_1,\ \ldots,\ C(f)^{n-1}e_1
\]
恰好构成一组基,因此 $e_1$ 是循环向量.
由最后一列可得
\[
C(f)^ne_1
+c_{n-1}C(f)^{n-1}e_1+\cdots+c_0e_1=0,
\]
即
\[
f(C(f))e_1=0.
\]
由于 $e_1$ 生成整个空间,
\[
f(C(f))=O.
\]
若存在次数小于 $n$ 的非零多项式 $g$ 使 $g(C(f))=O$,则
\[
e_1,C(f)e_1,\ldots,C(f)^{n-1}e_1
\]
线性相关,矛盾.因此最小多项式次数为 $n$,从而
\[
m_{C(f)}=f.
\]
最小多项式整除特征多项式,且二者均为 $n$ 次首一多项式,故
\[
\boxed{m_{C(f)}(\lambda)=f_{C(f)}(\lambda)=f(\lambda)}.
\]
\end{solution}


%例题6.1.3
\begin{example}
设 $\lambda_1,\lambda_2,\lambda_3$ 是
\[
f(x)=x^3+x^2+x+2
\]
的根,令
\[
h(x)=x^2+x+1.
\]
求一个有理系数多项式 $g(x)$,使
\[
h(\lambda_1),h(\lambda_2),h(\lambda_3)
\]
均为 $g(x)$ 的根.
\end{example}
\exampleblank

\begin{solution}
令
\[
y=h(x)=x^2+x+1.
\]
消去 $x$,计算结式
\[
\operatorname{Res}_x
\left(x^3+x^2+x+2,\ x^2+x+1-y\right)
=
-y^3+y^2-2y+4.
\]
故可取首一多项式
\[
\boxed{g(y)=y^3-y^2+2y-4}.
\]
\end{solution}


%例题6.1.4
\begin{example}[\markstar]
设 $A\in M_n(\mathbb{C})$,
\[
W=\{f(A):f(x)\in\mathbb{C}[x]\},
\]
且 $m_A(x)$ 为 $A$ 的最小多项式.证明
\[
\dim W=\deg m_A.
\]
\end{example}
\exampleblank

\begin{solution}
设
\[
m_A(x)=x^r+c_{r-1}x^{r-1}+\cdots+c_0.
\]
由带余除法,任意 $f(x)$ 都可写成
\[
f(x)=q(x)m_A(x)+r(x),
\qquad
\deg r<r.
\]
于是
\[
f(A)=r(A),
\]
故
\[
W=\Span\{I,A,\ldots,A^{r-1}\}.
\]
若这些矩阵线性相关,则存在次数小于 $r$ 的非零多项式零化 $A$,
与最小多项式的最小性矛盾.因此它们线性无关,
\[
\boxed{\dim W=r=\deg m_A}.
\]
\end{solution}


\subsection{准素分解}

\begin{lemma}[\markstar]
设 $f,g\in F[\lambda]$,$A$ 为有限维线性空间 $V$ 上的线性变换,则:
\begin{enumerate}
    \item $\ker I=\{0\}$,$\ker m_A(A)=V$;
    \item 若 $g\mid f$,则
    \[
    \ker g(A)\subseteq\ker f(A);
    \]
    \item 若 $d=(f,g)$,则
    \[
    \ker f(A)\cap\ker g(A)=\ker d(A);
    \]
    \item 若 $M=[f,g]$,则
    \[
    \ker f(A)+\ker g(A)=\ker M(A);
    \]
    \item 若 $(f,g)=1$,则
    \[
    \ker(fg)(A)=\ker f(A)\oplus\ker g(A).
    \]
\end{enumerate}
\end{lemma}


\begin{theorem}[\markstar\ 准素分解]
设
\[
m_A(\lambda)
=
p_1(\lambda)^{r_1}\cdots p_s(\lambda)^{r_s},
\]
其中 $p_i$ 是 $F[\lambda]$ 中两两不同的首一不可约多项式.令
\[
W_i=\ker p_i(A)^{r_i}.
\]
则:
\begin{enumerate}
    \item
    \[
    V=W_1\oplus\cdots\oplus W_s;
    \]
    \item $A_i=A|_{W_i}$ 的最小多项式为 $p_i^{r_i}$;
    \item 投影 $E_i:V\to W_i$ 可以表示成 $A$ 的多项式;
    \item 若 $B$ 与 $A$ 可交换,则每个 $W_i$ 都是 $B$ 的不变子空间.
\end{enumerate}
\end{theorem}

\begin{proposition}[\markstar]
若 $A$ 的特征多项式为
\[
f_A(\lambda)
=
p_1(\lambda)^{d_1}\cdots p_s(\lambda)^{d_s},
\]
则相应准素分量也可写为
\[
W_i
=
\ker p_i(A)^{d_i}
=
\bigcup_{k\geq1}\ker p_i(A)^k.
\]
当 $F=\mathbb{C}$ 时,$p_i(\lambda)=\lambda-\lambda_i$,
$W_i$ 称为属于 $\lambda_i$ 的根子空间.
\end{proposition}

\begin{proposition}[\marktwostar]
方阵 $A$ 在 $F$ 上可对角化,当且仅当其最小多项式在 $F$ 上分解为
两两互素的一次因子的乘积.
\end{proposition}

\begin{proposition}[\markstar]
若 $A$ 在 $F$ 中有 $n$ 个特征根(重根按重数计算),则存在唯一的
\[
A=D+N,
\qquad
DN=ND,
\]
其中 $D$ 在 $F$ 上可对角化,$N$ 为幂零矩阵,并且 $D,N$ 均为 $A$ 的多项式.
\end{proposition}


%例题6.1.5
\begin{example}
设
\[
A=
\begin{pmatrix}
0&-1&2&0\\
1&0&-2&0\\
0&0&1&0\\
1&1&-2&1
\end{pmatrix}
\]
定义 $\mathbb{R}^4$ 上的线性变换.求 $V$ 与 $A$ 的准素分解.
\end{example}
\exampleblank

\begin{solution}
特征多项式为
\[
f_A(\lambda)=(\lambda-1)^2(\lambda^2+1),
\]
而实际最小多项式为
\[
m_A(\lambda)=(\lambda-1)(\lambda^2+1).
\]
因此
\[
V=W_1\oplus W_2,
\qquad
W_1=\ker(A-I),\quad
W_2=\ker(A^2+I).
\]
可取
\[
W_1
=
\Span\left\{
\begin{pmatrix}2\\0\\1\\0\end{pmatrix},
\begin{pmatrix}0\\0\\0\\1\end{pmatrix}
\right\},
\]
\[
W_2
=
\Span\left\{
\begin{pmatrix}0\\1\\0\\0\end{pmatrix},
\begin{pmatrix}-1\\0\\0\\1\end{pmatrix}
\right\}.
\]
令 $P$ 的列依次取上述四个向量,则
\[
P^{-1}AP
=
\boxed{
\begin{pmatrix}
1&0&0&0\\
0&1&0&0\\
0&0&0&-1\\
0&0&1&0
\end{pmatrix}
}.
\]
这正是实数域上的准对角分解.
\end{solution}


%例题6.1.6
\begin{example}
设
\[
A=
\begin{pmatrix}
2&-1&2&2\\
1&0&-2&0\\
1&0&1&1\\
1&1&-2&-1
\end{pmatrix}.
\]
分别把它看作 $\mathbb{C}^4$ 与 $\mathbb{R}^4$ 上的线性变换,
求相应的准素分解.
\end{example}
\exampleblank

\begin{solution}
特征多项式为
\[
(\lambda-3)(\lambda+1)(\lambda^2+1).
\]
故在 $\mathbb{C}$ 上有四个一维根子空间,$A$ 可对角化,特征值为
\[
3,\ -1,\ i,\ -i.
\]

在 $\mathbb{R}$ 上,
\[
V=W_3\oplus W_{-1}\oplus W_q,
\]
其中
\[
W_3=\ker(A-3I)
=
\Span\left\{
\begin{pmatrix}2\\0\\1\\0\end{pmatrix}
\right\},
\]
\[
W_{-1}=\ker(A+I)
=
\Span\left\{
\begin{pmatrix}-2\\0\\-1\\4\end{pmatrix}
\right\},
\]
\[
W_q=\ker(A^2+I).
\]
可在 $W_q$ 中取
\[
v=
\begin{pmatrix}-1\\0\\2\\0\end{pmatrix},
\qquad
Av=
\begin{pmatrix}2\\-5\\1\\-5\end{pmatrix}.
\]
于是以
\[
P=
\left(
\begin{pmatrix}2\\0\\1\\0\end{pmatrix},
\begin{pmatrix}-2\\0\\-1\\4\end{pmatrix},
v,Av
\right)
\]
为过渡矩阵,
\[
\boxed{
P^{-1}AP
=
\begin{pmatrix}
3&0&0&0\\
0&-1&0&0\\
0&0&0&-1\\
0&0&1&0
\end{pmatrix}.
}
\]
\end{solution}


%例题6.1.7
\begin{example}
设
\[
V=W_1\oplus W_2\oplus\cdots\oplus W_s
\]
是 $V$ 关于线性变换 $A$ 的准素分解,而 $W$ 是 $A$ 的任意不变子空间.
证明
\[
W=(W_1\cap W)\oplus\cdots\oplus(W_s\cap W).
\]
\end{example}
\exampleblank

\begin{solution}
准素分解对应的正则投影 $E_i$ 都是 $A$ 的多项式.
因为 $W$ 对 $A$ 不变,所以也对任意 $A$ 的多项式不变,于是
\[
E_i(W)\subseteq W.
\]
对任意 $w\in W$,
\[
w=E_1w+\cdots+E_sw,
\]
且
\[
E_iw\in W_i\cap W.
\]
原准素分解本身是直和,因此
\[
\boxed{
W=(W_1\cap W)\oplus\cdots\oplus(W_s\cap W).
}
\]
\end{solution}


\subsection{循环分解}

\begin{definition}
对 $\alpha\in V$,定义
\[
F[A]\alpha
=
\{g(A)\alpha:g(\lambda)\in F[\lambda]\},
\]
称为由 $\alpha$ 生成的 $A$ 的\strongcn{循环子空间}.
若
\[
F[A]\alpha=V,
\]
则称 $V$ 为循环空间,$\alpha$ 为循环向量.

若 $g(A)\alpha=0$,则称 $g$ 为 $\alpha$ 的零化子;
次数最低的首一零化子称为 $\alpha$ 的最小零化子,记为 $m_\alpha$.
\end{definition}

\begin{theorem}[\markstar]
设
\[
W=F[A]\alpha.
\]
若
\[
\alpha,A\alpha,\ldots,A^{k-1}\alpha
\]
线性无关,而再加入 $A^k\alpha$ 后线性相关,并且
\[
A^k\alpha+c_{k-1}A^{k-1}\alpha+\cdots+c_0\alpha=0,
\]
则
\[
m_\alpha(\lambda)
=
\lambda^k+c_{k-1}\lambda^{k-1}+\cdots+c_0.
\]
此外:
\begin{enumerate}
    \item
    \[
    \dim W=k;
    \]
    \item
    \[
    \alpha,A\alpha,\ldots,A^{k-1}\alpha
    \]
    是 $W$ 的一组基;
    \item $A|_W$ 在该基下的矩阵为 $m_\alpha$ 的友矩阵;
    \item $A|_W$ 的最小多项式,特征多项式与 $m_\alpha$ 三者相等.
\end{enumerate}
\end{theorem}


%例题6.1.8
\begin{example}
设
\[
A=
\begin{pmatrix}
1&0&0\\
1&1&0\\
0&1&1
\end{pmatrix},
\qquad
\varepsilon_1=(1,0,0)^{\mathsf T}.
\]
求 $\varepsilon_1$ 的最小零化子,并给出相应循环基与矩阵表示.
\end{example}
\exampleblank

\begin{solution}
令
\[
N=A-I=
\begin{pmatrix}
0&0&0\\
1&0&0\\
0&1&0
\end{pmatrix}.
\]
则
\[
N\varepsilon_1=\varepsilon_2,\qquad
N^2\varepsilon_1=\varepsilon_3,\qquad
N^3\varepsilon_1=0.
\]
故
\[
\boxed{m_{\varepsilon_1}(\lambda)=(\lambda-1)^3}.
\]
并且
\[
\varepsilon_1,\ N\varepsilon_1,\ N^2\varepsilon_1
\]
是一组循环基,在此基下
\[
[A]
=
\boxed{
\begin{pmatrix}
1&0&0\\
1&1&0\\
0&1&1
\end{pmatrix}
}
\]
即一个转置约定下的 Jordan 链矩阵;反向排列基向量即可得到通常的
$J_3(1)$.
\end{solution}


%例题6.1.9
\begin{example}[\markstar]
设 $A$ 在某一基下的矩阵是首一多项式 $f$ 的友矩阵.
证明 $V$ 是循环空间.
\end{example}
\exampleblank

\begin{solution}
在友矩阵的标准基下,
\[
e_1,Ae_1,\ldots,A^{n-1}e_1
\]
恰好为一组基,因此
\[
F[A]e_1=V.
\]
故 $e_1$ 是循环向量.
\end{solution}


%例题6.1.10
\begin{example}
设
\[
A=\diag(3,3,-2)
\]
是 $\mathbb{R}^3$ 上的线性变换.
证明 $A$ 没有循环向量,并求
\[
v=(1,-1,3)^{\mathsf T}
\]
生成的循环子空间.
\end{example}
\exampleblank

\begin{solution}
$A$ 的最小多项式为
\[
m_A(\lambda)=(\lambda-3)(\lambda+2),
\]
次数只有 $2<3$.循环空间若等于整个 $\mathbb{R}^3$,其循环向量的最小零化子
必须等于 $A$ 的特征多项式并有次数 $3$,故不存在循环向量.

把
\[
v=(1,-1,0)^{\mathsf T}+(0,0,3)^{\mathsf T}
\]
分解到两个特征子空间,得到
\[
F[A]v
=
\boxed{
\Span\left\{
(1,-1,0)^{\mathsf T},\ (0,0,1)^{\mathsf T}
\right\}.
}
\]
\end{solution}


%例题6.1.11
\begin{example}
设
\[
A=
\begin{pmatrix}
1&i&0\\
-1&2&-i\\
0&1&1
\end{pmatrix}
\]
作用于 $\mathbb{C}^3$.求
\[
\alpha=(1,0,0)^{\mathsf T},
\qquad
\beta=(1,0,i)^{\mathsf T}
\]
的最小零化子.
\end{example}
\exampleblank

\begin{solution}
直接考察循环向量组
\[
\alpha,A\alpha,A^2\alpha,\ldots
\]
可得
\[
\boxed{
m_\alpha(x)
=
(x-1)(x^2-3x+2+2i)
}.
\]
而
\[
A\beta=\beta,
\]
所以
\[
\boxed{m_\beta(x)=x-1}.
\]
\end{solution}


%例题6.1.12
\begin{example}
设 $A$ 是 $n$ 维线性空间上的可对角化线性变换.证明:
\begin{enumerate}
    \item 若 $A$ 有循环向量,则 $A$ 有 $n$ 个互异特征值;
    \item 若 $A$ 有 $n$ 个互异特征值,$\alpha_1,\ldots,\alpha_n$
    为对应特征向量,则
    \[
    \alpha_1+\cdots+\alpha_n
    \]
    是循环向量.
\end{enumerate}
\end{example}
\exampleblank

\begin{solution}
若 $A$ 可对角化,则最小多项式中所有一次因子指数均为 $1$.
若又有循环向量,则
\[
m_A=f_A,
\]
所以特征多项式也无重根,即有 $n$ 个互异特征值.

反之设特征值 $\lambda_1,\ldots,\lambda_n$ 两两不同,令
\[
\alpha=\alpha_1+\cdots+\alpha_n.
\]
则
\[
\alpha,A\alpha,\ldots,A^{n-1}\alpha
\]
关于特征向量基的坐标矩阵是 Vandermonde 矩阵,
行列式为
\[
\displaystyle\prod_{i<j}(\lambda_j-\lambda_i)\neq0.
\]
故这些向量构成一组基,$\alpha$ 是循环向量.
\end{solution}


\begin{theorem}[\markstar\ 循环分解]
设 $A$ 是有限维线性空间 $V$ 上的线性变换,则存在非零向量
$\alpha_1,\ldots,\alpha_r$ 使
\[
V
=
F[A]\alpha_1\oplus\cdots\oplus F[A]\alpha_r.
\]
记各循环分量对应最小零化子为
\[
m_1,\ldots,m_r,
\]
可以安排成
\[
m_r\mid m_{r-1}\mid\cdots\mid m_1.
\]
这些多项式由 $A$ 唯一确定,称为 $A$ 的不变因子,并且
\[
m_A=m_1,
\qquad
f_A=m_1m_2\cdots m_r.
\]
\end{theorem}

\begin{theorem}[\markstar\ 有理标准形]
在上述记号下,存在可逆矩阵 $P$ 使
\[
P^{-1}AP
=
\diag\left(C(m_1),C(m_2),\ldots,C(m_r)\right).
\]
该块对角矩阵称为 $A$ 的有理标准形.
\end{theorem}


%例题6.1.13
\begin{example}
求下列实矩阵的有理标准形,并给出一组相应过渡矩阵:
\begin{enumerate}
    \item
    \[
    A_1=
    \begin{pmatrix}
    5&-6&-6\\
    -1&4&2\\
    3&-6&-4
    \end{pmatrix};
    \]
    \item
    \[
    A_2=
    \begin{pmatrix}
    3&-4&-4\\
    -1&3&2\\
    2&-4&-3
    \end{pmatrix};
    \]
    \item
    \[
    A_3=
    \begin{pmatrix}
    0&-1&2&0\\
    1&0&-2&0\\
    0&0&1&0\\
    1&1&-2&1
    \end{pmatrix}.
    \]
\end{enumerate}
\end{example}
\exampleblank

\begin{solution}
第一问的不变因子可取
\[
m_1=(\lambda-1)(\lambda-2),
\qquad
m_2=\lambda-2.
\]
因此有理标准形可写为
\[
B_1=
\diag\left(
2,\
\begin{pmatrix}0&-2\\1&3\end{pmatrix}
\right).
\]
取
\[
P_1=
\begin{pmatrix}
2&-2&-6\\
1&-2&-2\\
0&-2&2
\end{pmatrix},
\]
则
\[
\boxed{P_1^{-1}A_1P_1=B_1}.
\]

第二问
\[
f_{A_2}=(\lambda-1)^3,
\qquad
m_{A_2}=(\lambda-1)^2.
\]
可取不变因子
\[
m_1=(\lambda-1)^2,\qquad m_2=\lambda-1,
\]
有理标准形
\[
B_2=
\diag\left(
1,\
\begin{pmatrix}0&-1\\1&2\end{pmatrix}
\right).
\]
例如
\[
P_2=
\begin{pmatrix}
2&-2&-2\\
1&-2&-2\\
0&-2&0
\end{pmatrix}
\]
满足
\[
\boxed{P_2^{-1}A_2P_2=B_2}.
\]

第三问不变因子为
\[
m_1=(\lambda-1)(\lambda^2+1),
\qquad
m_2=\lambda-1.
\]
于是
\[
B_3=
\diag\left(
1,\
C\left((\lambda-1)(\lambda^2+1)\right)
\right),
\]
而
\[
(\lambda-1)(\lambda^2+1)=\lambda^3-\lambda^2+\lambda-1,
\]
故
\[
C=
\begin{pmatrix}
0&0&1\\
1&0&-1\\
0&1&1
\end{pmatrix}.
\]
取
\[
P_3=
\begin{pmatrix}
2&-2&-2&0\\
0&-2&0&-2\\
1&-2&-2&0\\
0&-1&0&1
\end{pmatrix},
\]
则
\[
\boxed{P_3^{-1}A_3P_3=B_3}.
\]
\end{solution}


%例题6.1.14
\begin{example}
设
\[
A=
\begin{pmatrix}
2&0&0&0\\
1&2&0&0\\
0&a&2&0\\
0&0&b&2
\end{pmatrix}.
\]
求特征多项式;并分别讨论
\[
(a,b)=(1,1),\qquad(0,0),\qquad(0,1)
\]
时的最小多项式与循环分解.
\end{example}
\exampleblank

\begin{solution}
无论 $a,b$ 如何,
\[
\boxed{f_A(\lambda)=(\lambda-2)^4}.
\]
记 $N=A-2I$.

当 $a=b=1$ 时,
\[
N^3e_1=e_4\neq0,\qquad N^4=0,
\]
故
\[
m_A=(\lambda-2)^4,
\]
且 $e_1$ 是循环向量:
\[
V=F[A]e_1.
\]

当 $a=b=0$ 时,$N^2=0$ 且 $N\neq0$,
\[
m_A=(\lambda-2)^2.
\]
可取循环分解
\[
V
=
F[A]e_1\oplus F[A]e_3\oplus F[A]e_4,
\]
对应块大小为 $2,1,1$.

当 $a=0,b=1$ 时,同样
\[
m_A=(\lambda-2)^2,
\]
可取
\[
\boxed{
V=F[A]e_1\oplus F[A]e_3,
}
\]
两个循环分量都为二维,对应两个 $J_2(2)$ 型块.
\end{solution}


%例题6.1.15
\begin{example}
设 $A$ 为 $F$ 上 $n$ 维线性空间 $V$ 的线性变换.
证明:$V$ 的每个非零向量都是循环向量,当且仅当
$A$ 的特征多项式在 $F$ 上不可约.
\end{example}
\exampleblank

\begin{solution}
若 $f_A$ 不可约,取任意 $0\neq v\in V$.其最小零化子 $m_v$ 整除
$m_A$,而
\[
m_A\mid f_A.
\]
由于 $f_A$ 不可约且 $m_v$ 非常数,只能有
\[
m_v=f_A.
\]
故
\[
\dim F[A]v=\deg m_v=n,
\]
所以 $v$ 是循环向量.

反之,若每个非零向量都是循环向量,而 $f_A$ 可约,则 $m_A=f_A$.
取 $f_A=gh$ 为两个正次数因子.由 Cayley--Hamilton,
\[
g(A)h(A)=0.
\]
若 $h(A)\neq0$,取 $v$ 使 $h(A)v\neq0$,则非零向量
$w=h(A)v$ 被 $g(A)$ 零化,其最小零化子次数小于 $n$,不可能循环.
若 $h(A)=0$,则 $h$ 本身已是次数小于 $n$ 的零化多项式,同样矛盾.
故 $f_A$ 必不可约.
\end{solution}


\begin{proposition}[\markstar]
对 $n$ 维空间上的线性变换 $A$,下列命题彼此等价:
\begin{enumerate}
    \item
    \[
    m_A=f_A;
    \]
    \item $V$ 是循环空间;
    \item 存在 $\alpha$ 使
    \[
    \alpha,A\alpha,\ldots,A^{n-1}\alpha
    \]
    为一组基;
    \item 与 $A$ 可交换的每个线性变换都是 $A$ 的多项式;
    \item 对每个特征值,其特征子空间都是一维;
    \item $A$ 在某一基下为一个友矩阵.
\end{enumerate}
\end{proposition}


%例题6.1.16
\begin{example}
设 $V$ 是有限维实线性空间,$T:V\to V$ 满足:
\begin{enumerate}
    \item $T$ 的最小多项式在 $\mathbb{R}$ 上不可约;
    \item 存在 $v\in V$ 使
    \[
    \{T^i(v):i\geq0\}
    \]
    张成 $V$.
\end{enumerate}
证明 $V$ 中不存在非平凡的 $T$-不变子空间.
\end{example}
\exampleblank

\begin{solution}
第二个条件说明 $V$ 是循环 $\,\mathbb{R}[x]\,$-模,
\[
V\cong \mathbb{R}[x]/(m_T).
\]
由于 $m_T$ 不可约,商环
\[
\mathbb{R}[x]/(m_T)
\]
是域.循环模的 $T$-不变子空间对应于该商环的理想,而域只有
$\{0\}$ 和自身两个理想.因此只有
\[
\boxed{\{0\},V}
\]
两个 $T$-不变子空间.
\end{solution}


%例题6.1.17
\begin{example}
设 $T$ 是 $n$ 维线性空间上的可对角化线性变换.证明下列命题等价:
\begin{enumerate}
    \item $T$ 在底域中有 $n$ 个互不相同的特征值;
    \item 存在 $\alpha\in V$ 使
    \[
    \alpha,T\alpha,\ldots,T^{n-1}\alpha
    \]
    线性无关.
\end{enumerate}
\end{example}
\exampleblank

\begin{solution}
这正是前面``可对角化变换存在循环向量''的等价表述.
若有 $n$ 个互异特征值,取各特征向量之和,利用 Vandermonde 行列式即可得到第二条.
反之,第二条说明 $\alpha$ 为循环向量,因此
\[
m_T=f_T.
\]
又因 $T$ 可对角化,$m_T$ 无重根,于是 $f_T$ 也无重根,
即有 $n$ 个互异特征值.
\end{solution}


\subsection{Jordan 标准形}

\begin{theorem}[\marktwostar\ Jordan 标准形]
对任一复方阵 $A$,存在可逆复矩阵 $P$ 使
\[
P^{-1}AP
=
\diag\{J_{k_{11}}(\lambda_1),\ldots,J_{k_{sr_s}}(\lambda_s)\}.
\]
其中各 $J_k(\lambda)$ 为 Jordan 块.
更一般地,只要 $A$ 的特征多项式在底域中完全分裂,同样结论成立.
\end{theorem}

\begin{proposition}[\markstar]
属于特征值 $\lambda$ 的 Jordan 块个数等于特征子空间
\[
\ker(A-\lambda I)
\]
的维数.
\end{proposition}


%例题6.1.18
\begin{example}[\markstar]
设 $J_{10}(0)$ 为 $10$ 阶幂零 Jordan 块,求
\[
J_{10}(0)^3
\]
的 Jordan 标准形.
\end{example}
\exampleblank

\begin{solution}
对 $J_n(0)^k$,若
\[
n=qk+r,\qquad0\leq r<k,
\]
则其 Jordan 块为 $r$ 个 $J_{q+1}(0)$ 与 $k-r$ 个 $J_q(0)$.
这里
\[
10=3\cdot3+1,
\]
故
\[
\boxed{
J_{10}(0)^3
\sim
J_4(0)\oplus J_3(0)\oplus J_3(0).
}
\]
\end{solution}


%例题6.1.19
\begin{example}[\markstar]
设 $J_n(a)$ 为 $n$ 阶 Jordan 块,求
\[
J_n(a)^2
\]
的 Jordan 标准形.
\end{example}
\exampleblank

\begin{solution}
若 $a\neq0$,函数 $f(x)=x^2$ 在 $a$ 处满足
\[
f'(a)=2a\neq0,
\]
故 Jordan 块大小保持不变:
\[
\boxed{J_n(a)^2\sim J_n(a^2)}.
\]

若 $a=0$,写
\[
n=2q+r,\qquad r\in\{0,1\},
\]
则
\[
\boxed{
J_n(0)^2
\sim
\begin{cases}
J_q(0)\oplus J_q(0),&r=0,\\[2mm]
J_{q+1}(0)\oplus J_q(0),&r=1.
\end{cases}
}
\]
\end{solution}


\begin{proposition}[\marktwostar]
更一般地,对 $J_n(a)^k$:
\begin{enumerate}
    \item 若 $a\neq0$,则
    \[
    J_n(a)^k\sim J_n(a^k);
    \]
    \item 若 $a=0$,写
    \[
    n=qk+r,\qquad0\leq r<k,
    \]
    则
    \[
    J_n(0)^k
    \sim
    \underbrace{J_{q+1}(0)\oplus\cdots\oplus J_{q+1}(0)}_{r\text{ 个}}
    \oplus
    \underbrace{J_q(0)\oplus\cdots\oplus J_q(0)}_{k-r\text{ 个}},
    \]
    其中零阶块应忽略.
\end{enumerate}
\end{proposition}

\begin{proposition}[\markstar]
设属于 $\lambda_0$ 的 $k$ 阶 Jordan 块个数为 $\delta_k$,则
\[
\boxed{
\delta_k
=
\rank(\lambda_0I-A)^{k-1}
+\rank(\lambda_0I-A)^{k+1}
-2\rank(\lambda_0I-A)^k.
}
\]
\end{proposition}


\newpage

\section{三种因子(Jordan 标准型 2)}

\subsection{三种因子的定义与性质}

\begin{definition}
对多项式矩阵 $A(\lambda)$,以下操作称为初等行变换:
交换两行;某行乘以非零常数;某行乘以多项式后加到另一行.
列变换同理.
\end{definition}

\begin{theorem}[\markstar\ Smith 标准形]
设 $A(\lambda)$ 是多项式矩阵,则存在可逆多项式矩阵
$P(\lambda),Q(\lambda)$,使
\[
P(\lambda)A(\lambda)Q(\lambda)
=
\diag\left(d_1(\lambda),\ldots,d_r(\lambda),0,\ldots,0\right),
\]
其中
\[
d_1\mid d_2\mid\cdots\mid d_r.
\]
$d_i$ 称为 $A(\lambda)$ 的不变因子.
\end{theorem}

\begin{definition}
$A(\lambda)$ 的全部 $k$ 阶非零子式的首一最大公因式记为
$D_k(\lambda)$,称为第 $k$ 阶行列式因子.它们满足
\[
D_k=d_1d_2\cdots d_k,
\]
故
\[
d_1=D_1,\qquad
d_k=\dfrac{D_k}{D_{k-1}}.
\]
\end{definition}

\begin{definition}
把所有非平凡不变因子在底域上分解为不可约多项式的幂,
所得全部准素因子(重复者重复计入)称为初等因子.
\end{definition}

\begin{proposition}[\markstar]
块对角多项式矩阵
\[
\diag(A_1(\lambda),A_2(\lambda))
\]
的初等因子组等于两个块的初等因子组的合并.
\end{proposition}

\begin{theorem}[\marktwostar]
对同阶方阵 $A,B$,下列命题等价:
\begin{enumerate}
    \item $A\sim B$;
    \item $\lambda I-A$ 与 $\lambda I-B$ 相抵;
    \item 两者不变因子组相同;
    \item 两者行列式因子组相同;
    \item 两者初等因子组相同.
\end{enumerate}
\end{theorem}

\begin{theorem}[\markstar]
若复矩阵 $A$ 的初等因子为
\[
(\lambda-\lambda_1)^{k_1},\ldots,
(\lambda-\lambda_s)^{k_s},
\]
则
\[
A
\sim
\diag\left(
J_{k_1}(\lambda_1),\ldots,J_{k_s}(\lambda_s)
\right).
\]
\end{theorem}


%例题6.2.1
\begin{example}
求下列多项式矩阵的 Smith 标准形:
\begin{enumerate}
    \item
    \[
    \begin{pmatrix}
    \lambda&1\\
    0&\lambda
    \end{pmatrix};
    \]
    \item
    \[
    \begin{pmatrix}
    \lambda^2-1&\lambda+1\\
    \lambda+1&\lambda^2+2\lambda+1
    \end{pmatrix};
    \]
    \item
    \[
    \begin{pmatrix}
    \lambda^2-1&0\\
    0&(\lambda-1)^3
    \end{pmatrix};
    \]
    \item
    \[
    \begin{pmatrix}
    \lambda+1&\lambda^2+1&\lambda^2\\
    3\lambda-1&3\lambda^2-1&\lambda^2+2\lambda\\
    \lambda-1&\lambda^2-1&\lambda
    \end{pmatrix}.
    \]
\end{enumerate}
\end{example}
\exampleblank

\begin{solution}
第一问所有 $1$ 阶子式的最大公因式为 $1$,行列式为 $\lambda^2$,故
\[
\boxed{\diag(1,\lambda^2)}.
\]

第二问
\[
D_1=\lambda+1,
\]
且
\[
\det A(\lambda)
=
(\lambda+1)^2(\lambda^2-2).
\]
故
\[
d_2
=
\dfrac{D_2}{D_1}
=
(\lambda+1)(\lambda^2-2),
\]
所以
\[
\boxed{
\diag\left(
\lambda+1,\,
(\lambda+1)(\lambda^2-2)
\right).
}
\]

第三问
\[
D_1=\lambda-1,
\]
\[
D_2=(\lambda^2-1)(\lambda-1)^3
=(\lambda-1)^4(\lambda+1),
\]
故
\[
\boxed{
\diag\left(
\lambda-1,\,
(\lambda-1)^3(\lambda+1)
\right).
}
\]

第四问原矩阵秩为 $2$,$D_1=1$,全部二阶子式的最大公因式为 $\lambda$,
故 Smith 标准形为
\[
\boxed{\diag(1,\lambda,0)}.
\]
\end{solution}


%例题6.2.2
\begin{example}
求下列方阵在 $\mathbb{C}$ 上的 Jordan 标准形,并给出可逆矩阵 $P$
使 $P^{-1}AP=J$:
\begin{enumerate}
    \item
    \[
    \begin{pmatrix}
    2&1&3\\
    0&2&-1\\
    0&0&2
    \end{pmatrix};
    \]
    \item
    \[
    \begin{pmatrix}
    2&-1&-1\\
    2&-1&-2\\
    -1&1&2
    \end{pmatrix};
    \]
    \item
    \[
    \begin{pmatrix}
    4&-5&2\\
    5&-7&3\\
    6&-9&4
    \end{pmatrix};
    \]
    \item
    \[
    \begin{pmatrix}
    1&-3&0&3\\
    -2&-6&0&13\\
    0&-3&1&3\\
    -1&-4&0&8
    \end{pmatrix};
    \]
    \item
    \[
    \begin{pmatrix}
    3&-1&0&0\\
    1&1&0&0\\
    3&0&5&-3\\
    4&-1&3&-1
    \end{pmatrix}.
    \]
\end{enumerate}
\end{example}
\exampleblank

\begin{solution}
可分别取:

\[
J_1=J_3(2),\qquad
P_1=
\begin{pmatrix}
-1&3&0\\
0&-1&0\\
0&0&1
\end{pmatrix}.
\]

\[
J_2=J_2(1)\oplus(1),\qquad
P_2=
\begin{pmatrix}
1&1&1\\
2&0&1\\
-1&0&0
\end{pmatrix}.
\]

\[
J_3=J_2(0)\oplus(1),\qquad
P_3=
\begin{pmatrix}
-1&1&1\\
-2&1&1\\
-3&0&1
\end{pmatrix}.
\]

\[
J_4=J_3(1)\oplus(1),\qquad
P_4=
\begin{pmatrix}
3&0&1&0\\
1&-2&0&0\\
3&0&0&1\\
1&-1&0&0
\end{pmatrix}.
\]

\[
J_5=J_2(2)\oplus J_2(2),\qquad
P_5=
\begin{pmatrix}
1&1&-1&0\\
1&0&-1&1\\
3&0&0&0\\
4&0&-1&0
\end{pmatrix}.
\]

直接相乘可验证
\[
\boxed{P_i^{-1}A_iP_i=J_i}.
\]
\end{solution}


%例题6.2.3
\begin{example}
求下列矩阵的不变因子与初等因子:
\begin{enumerate}
    \item 原讲义给出的一个含参数 $a,b$ 的 $6$ 阶矩阵;
    \item
    \[
    L(\lambda)=
    \begin{pmatrix}
    \lambda&0&\cdots&0&a_n\\
    -1&\lambda&\cdots&0&a_{n-1}\\
    0&-1&\ddots&0&a_{n-2}\\
    \vdots&\ddots&\ddots&\lambda&\vdots\\
    0&\cdots&0&-1&\lambda+a_1
    \end{pmatrix};
    \]
    \item $A=\alpha^{\mathsf T}\beta$,且
    \[
    \operatorname{tr}A=0;
    \]
    \item 上三角矩阵
    \[
    A=
    \begin{pmatrix}
    a&b&\cdots&b\\
    0&a&\ddots&\vdots\\
    \vdots&\ddots&\ddots&b\\
    0&\cdots&0&a
    \end{pmatrix}.
    \]
\end{enumerate}
\end{example}
\exampleblank

\begin{solution}
第一问中原 PDF 只给出含参数矩阵而没有说明 $a,b$ 的取值条件.
不变因子会在 $a=0$,$b=0$,$a=\pm b$ 等退化情形发生改变,因此不能在
不分类的前提下给出唯一答案.本题已列入《解答问题记录》,正式定稿时应先明确参数条件.

第二问正是某个 $n$ 次首一多项式
\[
p(\lambda)
=
\lambda^n+a_1\lambda^{n-1}+\cdots+a_n
\]
的友矩阵对应的 $\lambda I-C(p)$ 型矩阵.
其前 $n-1$ 个 Smith 因子均为 $1$,最后一个为 $p(\lambda)$:
\[
\boxed{
1,\ldots,1,p(\lambda).
}
\]
初等因子即 $p$ 在底域上的准素因子.

第三问若 $A\neq0$,则 $r(A)=1$ 且
\[
A^2=(\operatorname{tr}A)A=0.
\]
故 Jordan 型为
\[
J_2(0)\oplus 0\oplus\cdots\oplus0.
\]
因此初等因子为
\[
\boxed{\lambda^2,\ \underbrace{\lambda,\ldots,\lambda}_{n-2\text{ 个}}},
\]
相应不变因子可写为
\[
\boxed{
1,\ 
\underbrace{\lambda,\ldots,\lambda}_{n-2\text{ 个}},
\lambda^2.
}
\]

第四问若 $b\neq0$,令 $N=A-aI$,则
\[
N^{n-1}\neq0,\qquad N^n=0,
\]
所以 $A$ 相似于单个 Jordan 块 $J_n(a)$.
故
\[
\boxed{
1,\ldots,1,(\lambda-a)^n
}
\]
是不变因子组的非平凡部分,初等因子仅有
\[
\boxed{(\lambda-a)^n}.
\]
若 $b=0$,则 $A=aI$,应另行退化为 $n$ 个一阶块.
\end{solution}


%例题6.2.4
\begin{example}
设 $A$ 为 $n$ 阶方阵,
\[
|A|=18,
\qquad
3A+A^*=15I_n,
\]
其中 $A^*$ 为伴随矩阵.求 $A$ 的最小多项式与 Jordan 标准形.
\end{example}
\exampleblank

\begin{solution}
因为 $\det A=18\neq0$,
\[
A^*=18A^{-1}.
\]
原式右乘 $A$:
\[
3A^2-15A+18I=0,
\]
即
\[
A^2-5A+6I=0,
\]
所以
\[
m_A(\lambda)\mid(\lambda-2)(\lambda-3).
\]
$A$ 可对角化,特征值只能为 $2,3$.设重数为 $r,s$,则
\[
2^r3^s=18=2\cdot3^2,
\]
故
\[
r=1,\qquad s=2,\qquad n=3.
\]
于是
\[
\boxed{
m_A(\lambda)=(\lambda-2)(\lambda-3),
\qquad
J=\diag(2,3,3).
}
\]
\end{solution}


%例题6.2.5
\begin{example}
设 $A$ 为秩 $n-1$ 的 $n$ 阶矩阵,特征值为
\[
\lambda_1,\ldots,\lambda_n,
\qquad
\lambda_n=0.
\]
求 $A^*$ 的 Jordan 标准形.
\end{example}
\exampleblank

\begin{solution}
由于 $r(A)=n-1$,
\[
r(A^*)=1.
\]
又 $AA^*=0$,故
\[
\operatorname{Im}A^*\subseteq\ker A.
\]

若
\[
\mu=\lambda_1\lambda_2\cdots\lambda_{n-1}\neq0,
\]
则零特征值为单根,$A^*$ 的唯一非零特征值为 $\mu$,于是
\[
\boxed{
A^*\sim\diag(\mu,0,\ldots,0).
}
\]

若 $\mu=0$,则 $A^*$ 的全部特征值都为 $0$,但秩为 $1$,故
\[
(A^*)^2=0,\qquad A^*\neq0,
\]
从而
\[
\boxed{
A^*\sim J_2(0)\oplus0\oplus\cdots\oplus0.
}
\]
\end{solution}


%例题6.2.6
\begin{example}
判断下列两组矩阵是否相似:
\begin{enumerate}
    \item
    \[
    A=
    \begin{pmatrix}
    1&0&0&0\\
    0&-1&0&0\\
    0&0&0&1\\
    0&0&0&0
    \end{pmatrix},
    \quad
    B=
    \begin{pmatrix}
    -1&0&0&0\\
    -1&1&1&-1\\
    -1&0&0&0\\
    -1&0&1&0
    \end{pmatrix};
    \]
    \item
    \[
    A=
    \begin{pmatrix}
    4&6&-15\\
    1&3&-5\\
    1&2&-4
    \end{pmatrix},
    \quad
    B=
    \begin{pmatrix}
    1&-3&3\\
    -2&-6&13\\
    -1&-4&8
    \end{pmatrix}.
    \]
\end{enumerate}
\end{example}
\exampleblank

\begin{solution}
第一组的 Jordan 标准形均为
\[
\boxed{
(-1)\oplus J_2(0)\oplus(1),
}
\]
故二者相似.

第二组的特征多项式均为
\[
(\lambda-1)^3,
\]
但第一矩阵的 Jordan 型为
\[
J_2(1)\oplus(1),
\]
第二矩阵的 Jordan 型为
\[
J_3(1).
\]
因此
\[
\boxed{\text{第二组不相似}.}
\]
\end{solution}


%例题6.2.7
\begin{example}
设 $A,B,C,D\in M_3(\mathbb{C})$ 具有相同的特征多项式.
证明其中必有两个矩阵相似.
\end{example}
\exampleblank

\begin{solution}
三阶复矩阵在特征多项式固定后,Jordan 型最多只有三种:
\begin{enumerate}
    \item 三个互异特征值时只有一种;
    \item 一个二重根与一个单根时,对二重根只有
    $J_2$ 或两个 $J_1$ 两种;
    \item 只有一个三重特征值时,Jordan 块大小只有分拆
    \[
    3,\qquad2+1,\qquad1+1+1
    \]
    三种.
\end{enumerate}
因此相似类最多三个.四个矩阵由抽屉原理必有两个属于同一相似类.
\end{solution}


%例题6.2.8
\begin{example}
设 $A,B\in M_3(\mathbb{C})$ 都只有一个特征值 $\lambda_0$.
证明
\[
A\sim B
\iff
\dim V_{\lambda_0}(A)=\dim V_{\lambda_0}(B).
\]
\end{example}
\exampleblank

\begin{solution}
三阶,单一特征值的 Jordan 型只有
\[
J_3(\lambda_0),\qquad
J_2(\lambda_0)\oplus J_1(\lambda_0),\qquad
\lambda_0I_3.
\]
它们的特征子空间维数依次为
\[
1,\qquad2,\qquad3.
\]
因此该维数唯一决定 Jordan 型,故结论成立.
\end{solution}


%例题6.2.9
\begin{example}[\markstar]
设 $A$ 为 $n$ 阶复方阵,全部特征值均为 $1$.
证明对任意正整数 $k$,
\[
A^k\sim A.
\]
\end{example}
\exampleblank

\begin{solution}
把 $A$ 化为 Jordan 标准形.对任意块
\[
J_s(1)=I+N,
\]
有
\[
J_s(1)^k
=
I+kN+\cdots.
\]
因为
\[
k\neq0
\]
在复数域中为非零数,所以一次项系数不为零,Jordan 块大小保持不变.
因此
\[
J_s(1)^k\sim J_s(1).
\]
逐块合并得到
\[
\boxed{A^k\sim A}.
\]
\end{solution}


%例题6.2.10
\begin{example}[\markstar]
设 $A$ 是复方阵,全部特征值均为 $\pm1$.证明
\[
A^{\mathsf T}\sim A^{-1}.
\]
\end{example}
\exampleblank

\begin{solution}
任意矩阵与其转置相似,所以
\[
A^{\mathsf T}\sim A.
\]
另一方面,对 Jordan 块 $J_s(\lambda)$,$\lambda=\pm1$,
函数
\[
f(x)=x^{-1}
\]
在 $\lambda$ 处满足
\[
f'(\lambda)=-\lambda^{-2}\neq0,
\]
故
\[
J_s(\lambda)^{-1}\sim J_s(\lambda^{-1})=J_s(\lambda).
\]
因此
\[
A^{-1}\sim A.
\]
于是
\[
\boxed{A^{\mathsf T}\sim A^{-1}}.
\]
\end{solution}


\subsection{平方根问题}

\begin{proposition}[\markstar]
任一可逆复矩阵 $A$ 都存在平方根,即存在复矩阵 $B$ 使
\[
B^2=A.
\]
\end{proposition}


%例题6.2.11
\begin{example}
设
\[
A=
\begin{pmatrix}
0&2011&1\\
0&0&2011\\
0&0&0
\end{pmatrix}.
\]
证明方程
\[
X^2=A
\]
在三阶复方阵中无解.
\end{example}
\exampleblank

\begin{solution}
$A$ 为幂零矩阵,且
\[
A^2=
\begin{pmatrix}
0&0&2011^2\\
0&0&0\\
0&0&0
\end{pmatrix}
\neq0,
\qquad
A^3=0.
\]
若 $X^2=A$,则 $X$ 的所有特征值平方均为 $0$,故 $X$ 也是三阶幂零矩阵,
所以
\[
X^3=0.
\]
于是
\[
A^2=X^4=0,
\]
矛盾.
\end{solution}


%例题6.2.12
\begin{example}
设
\[
A=
\begin{pmatrix}
-2&1&0\\
-4&2&0\\
-2&1&0
\end{pmatrix}.
\]
判断 $A$ 是否有平方根;若有,给出一个.
\end{example}
\exampleblank

\begin{solution}
有
\[
A^2=0,\qquad r(A)=1,
\]
故
\[
A\sim J_2(0)\oplus(0).
\]
而
\[
J_3(0)^2\sim J_2(0)\oplus(0),
\]
所以 $A$ 有平方根.

例如可取
\[
\boxed{
B=
\begin{pmatrix}
\dfrac12&-\dfrac94&4\\[2mm]
1&-\dfrac92&8\\
0&-2&4
\end{pmatrix}.
}
\]
直接计算可得
\[
B^2=A.
\]
\end{solution}


%例题6.2.13
\begin{example}
设 $A$ 是 $F$ 上的 $n$ 阶幂零矩阵,且幂零指数为 $n$:
\[
A^{n-1}\neq0,\qquad A^n=0.
\]
证明不存在 $n$ 阶方阵 $B$ 使
\[
A=B^2.
\]
\end{example}
\exampleblank

\begin{solution}
若 $A=B^2$,则 $B$ 也只有零特征值,因而是 $n$ 阶幂零矩阵,
所以
\[
B^n=0.
\]
令
\[
m=\left\lceil\dfrac n2\right\rceil.
\]
则 $m\leq n-1$($n\geq2$),并且
\[
A^m=B^{2m}=0
\]
因为 $2m\geq n$.这与 $A$ 的幂零指数恰为 $n$ 矛盾.
\end{solution}


\subsection{\texorpdfstring{$A^n$}{A 的 n 次方} 问题}

原讲义把常见方法分为数学归纳法,秩一矩阵法,零化多项式法与相似标准形法.

%例题6.2.14
\begin{example}
设
\[
A=
\begin{pmatrix}
\cos\varphi&-\sin\varphi\\
\sin\varphi&\cos\varphi
\end{pmatrix}.
\]
求 $A^n$.
\end{example}
\exampleblank

\begin{solution}
$A$ 表示旋转角 $\varphi$,连续复合 $n$ 次即旋转 $n\varphi$:
\[
\boxed{
A^n=
\begin{pmatrix}
\cos n\varphi&-\sin n\varphi\\
\sin n\varphi&\cos n\varphi
\end{pmatrix}.
}
\]
\end{solution}


%例题6.2.15
\begin{example}
设
\[
A=
\begin{pmatrix}
a&c\\
0&b
\end{pmatrix}.
\]
证明
\[
A^n=
\begin{pmatrix}
a^n&
\left(a^{n-1}+a^{n-2}b+\cdots+b^{n-1}\right)c\\
0&b^n
\end{pmatrix}.
\]
\end{example}
\exampleblank

\begin{solution}
对 $n$ 归纳即可.若结论对 $n$ 成立,则右乘 $A$ 后右上角变为
\[
c\left(
a^n+a^{n-1}b+\cdots+b^n
\right),
\]
恰为 $n+1$ 时公式.
\end{solution}


\begin{proposition}[\markstar]
若 $R$ 为秩一矩阵且
\[
\operatorname{tr}R=k,
\]
则
\[
R^2=kR.
\]
因此:
\begin{enumerate}
    \item 若 $k=0$,
    \[
    (aI+R)^n=a^nI+na^{n-1}R;
    \]
    \item 若 $k\neq0$,
    \[
    (aI+R)^n
    =
    a^nI
    +
    \dfrac{(a+k)^n-a^n}{k}R.
    \]
\end{enumerate}
\end{proposition}


%例题6.2.16
\begin{example}
原讲义在``秩一矩阵''方法下给出
\[
A=
\begin{pmatrix}
5&2&4\\
-2&0&-8\\
3&6&16
\end{pmatrix}
\]
并写``求 $A$''.结合本节标题与上下文,应理解为求 $A^n$.
\end{example}
\exampleblank

\begin{solution}
特征多项式为
\[
(\lambda-4)^2(\lambda-13).
\]
令
\[
R=A-4I.
\]
则
\[
r(R)=1,\qquad \operatorname{tr}R=9.
\]
应用上一命题:
\[
\boxed{
A^n
=
4^nI
+
\dfrac{13^n-4^n}{9}(A-4I).
}
\]
原题文字``求 $A$''显然与本节内容不符,已在问题记录中注明.
\end{solution}


%例题6.2.17
\begin{example}
设 $A$ 为 $n$ 阶矩阵,对角元均为 $2$,非对角元均为 $1$.
求 $A^m$.
\end{example}
\exampleblank

\begin{solution}
写成
\[
A=I+J,
\]
其中 $J$ 为全 $1$ 矩阵.$J$ 为秩一矩阵且
\[
J^2=nJ.
\]
故
\[
\boxed{
A^m
=
I+\dfrac{(n+1)^m-1}{n}J.
}
\]
\end{solution}


%例题6.2.18
\begin{example}
设
\[
A=
\begin{pmatrix}
1&4&2\\
0&-3&4\\
0&4&3
\end{pmatrix}.
\]
求 $A^n$.
\end{example}
\exampleblank

\begin{solution}
特征值为
\[
-5,\quad1,\quad5
\]
且两两不同,因此最小多项式为
\[
(\lambda+5)(\lambda-1)(\lambda-5).
\]
将 $\lambda^n$ 对该三次多项式取余,设
\[
\lambda^n\equiv c_0+c_1\lambda+c_2\lambda^2,
\]
由在 $-5,1,5$ 三点插值得
\[
c_0=
\dfrac{(-5)^n}{12}
-\dfrac{5^n}{8}
+\dfrac{25}{24},
\]
\[
c_1=
\dfrac{5^n-(-5)^n}{10},
\]
\[
c_2=
\dfrac{(-5)^n}{60}
+\dfrac{5^n}{40}
-\dfrac1{24}.
\]
故
\[
\boxed{
A^n=c_0I+c_1A+c_2A^2.
}
\]
\end{solution}


%例题6.2.19
\begin{example}
设
\[
A=
\begin{pmatrix}
2&1&0\\
0&2&1\\
0&0&2
\end{pmatrix},
\qquad
f(x)=1+x+x^2+\cdots+x^7.
\]
求 $f(A)$.
\end{example}
\exampleblank

\begin{solution}
写
\[
A=2I+N,
\qquad
N^3=0.
\]
由 Taylor 展开,
\[
f(A)=f(2)I+f'(2)N+\dfrac{f''(2)}{2}N^2.
\]
计算得
\[
f(2)=255,\qquad
f'(2)=769,\qquad
\dfrac{f''(2)}2=1023.
\]
故
\[
\boxed{
f(A)=
255I+769N+1023N^2
=
\begin{pmatrix}
255&769&1023\\
0&255&769\\
0&0&255
\end{pmatrix}.
}
\]
\end{solution}


%例题6.2.20
\begin{example}
求下列矩阵的 $n$ 次方:
\begin{enumerate}
    \item
    \[
    A_1=
    \begin{pmatrix}
    0&\dfrac12&\dfrac12\\
    1&-\dfrac12&\dfrac12\\
    1&-\dfrac12&\dfrac12
    \end{pmatrix};
    \]
    \item
    \[
    A_2=
    \begin{pmatrix}
    2&1&3\\
    0&2&-1\\
    0&0&2
    \end{pmatrix};
    \]
    \item
    \[
    A_3=
    \begin{pmatrix}
    2&-1&-1\\
    2&-1&-2\\
    -1&1&2
    \end{pmatrix}.
    \]
\end{enumerate}
\end{example}
\exampleblank

\begin{solution}
第一问的特征值为 $-1,0,1$,且可对角化.因此对 $n\geq1$,
\[
\boxed{
A_1^n=
\begin{cases}
A_1,&n\text{ 为奇数},\\
A_1^2,&n\text{ 为偶数}.
\end{cases}
}
\]

第二问写
\[
A_2=2I+N,
\qquad
N=
\begin{pmatrix}
0&1&3\\
0&0&-1\\
0&0&0
\end{pmatrix},
\qquad N^3=0.
\]
故
\[
\boxed{
A_2^n
=
2^nI
+n2^{n-1}N
+C_2^n2^{n-2}N^2.
}
\]

第三问写
\[
A_3=I+N,
\qquad N^2=0.
\]
于是
\[
\boxed{
A_3^n=I+n(A_3-I).
}
\]
\end{solution}

\begin{remark}[快速处理矩阵幂]
原讲义给出的顺序是:先检查幂等,对合,幂零,秩一以及便于归纳的结构;
其次观察是否为数量矩阵加秩一矩阵;再考虑较低次零化多项式;
最后使用 Jordan 标准形或其他相似标准形作为通法.
\end{remark}


\newpage

\section{可对角化问题深化}

\begin{proposition}[\marktwostar]
设 $\varphi$ 是复数域上 $n$ 维线性空间 $V$ 的线性变换,$A$ 为其矩阵.
下列命题等价:
\begin{enumerate}
    \item $A$ 可对角化;
    \item $A$ 有 $n$ 个线性无关的特征向量;
    \item $A$ 的全部初等因子均为一次多项式;
    \item $A$ 的全部不变因子均无重根;
    \item $A$ 的最小多项式无重根;
    \item 对任意 $c\in\mathbb{C}$,
    \[
    \rank(cI-A)=\rank(cI-A)^2;
    \]
    \item 对任意特征值 $\lambda_0$,
    \[
    \operatorname{Im}(\varphi-\lambda_0I)
    \cap
    \ker(\varphi-\lambda_0I)
    =
    \{0\}.
    \]
\end{enumerate}
\end{proposition}


%例题6.3.1
\begin{example}[\markstar]
设 $\varphi$ 是 $n$ 维复线性空间 $V$ 的线性变换,
$\lambda_0$ 是特征值,$m_0$ 是 $\lambda_0$ 在最小多项式中的重数.
证明
\[
m_0
=
\min
\left\{
k\in\mathbb{Z}_+:
\ker(\lambda_0I-\varphi)^k
=
\ker(\lambda_0I-\varphi)^{k+1}
\right\}.
\]
\end{example}
\exampleblank

\begin{solution}
限制到 $\lambda_0$ 的根子空间 $W$,令
\[
N=(\varphi-\lambda_0I)|_W.
\]
则 $N$ 为幂零变换,其幂零指数恰等于最小多项式中
$(x-\lambda_0)$ 的指数 $m_0$.对幂零变换,
\[
\ker N\subsetneq\ker N^2\subsetneq\cdots\subsetneq\ker N^{m_0}=W,
\]
而当 $k\geq m_0$ 时核不再增大.因此第一次稳定正发生在 $k=m_0$.
\end{solution}


%例题6.3.2
\begin{example}[\markstar]
判断下列矩阵是否可对角化:
\begin{enumerate}
    \item 循环矩阵
    \[
    A=
    \begin{pmatrix}
    a_0&a_1&\cdots&a_{n-1}\\
    a_{n-1}&a_0&\cdots&a_{n-2}\\
    \vdots&\vdots&&\vdots\\
    a_1&a_2&\cdots&a_0
    \end{pmatrix};
    \]
    \item 幂等矩阵 $A^2=A$;
    \item 对合矩阵 $A^2=I$;
    \item
    \[
    A=
    \begin{pmatrix}
    I_n&I_n\\
    I_n&0
    \end{pmatrix}.
    \]
\end{enumerate}
\end{example}
\exampleblank

\begin{solution}
第一问在 $\mathbb{C}$ 上总可对角化.离散 Fourier 基
\[
(1,\omega,\omega^2,\ldots,\omega^{n-1})^{\mathsf T}
\]
给出 $n$ 个线性无关特征向量.

第二问最小多项式整除
\[
x(x-1),
\]
无重根,因此可对角化.

第三问最小多项式整除
\[
(x-1)(x+1).
\]
若底域特征不为 $2$,则可对角化;在特征 $2$ 的域上该结论需另论.

第四问直接计算
\[
A^2-A-I_{2n}=0.
\]
最小多项式整除
\[
x^2-x-1,
\]
该多项式在 $\mathbb{C}$(以及 $\mathbb{R}$)上有两个不同的根,
故
\[
\boxed{A\text{ 可对角化}.}
\]
\end{solution}


%例题6.3.3
\begin{example}
设 $A$ 是线性空间上的线性变换,且
\[
\rank A=1.
\]
证明:$A$ 或可对角化,或为幂零变换.
\end{example}
\exampleblank

\begin{solution}
秩一矩阵满足
\[
A^2=(\operatorname{tr}A)A.
\]
若
\[
\operatorname{tr}A\neq0,
\]
则最小多项式整除
\[
x(x-\operatorname{tr}A),
\]
其根互异,故 $A$ 可对角化.
若
\[
\operatorname{tr}A=0,
\]
则
\[
A^2=0,
\]
所以 $A$ 为幂零变换.
\end{solution}


%例题6.3.4
\begin{example}
设 $A$ 为 $n$ 阶实方阵,把同一矩阵分别看作
$\mathbb{R}^n$ 与 $\mathbb{C}^n$ 上的线性变换.
若 $A$ 在 $\mathbb{R}^n$ 中只有
\[
\{0\},\ \mathbb{R}^n
\]
两个不变子空间,证明它在 $\mathbb{C}^n$ 上可对角化.
\end{example}
\exampleblank

\begin{solution}
若 $n>2$,实特征多项式必有一个次数为 $1$ 或 $2$ 的不可约因子,
对应的核空间给出维数 $1$ 或 $2$ 的非平凡实不变子空间,矛盾.
因此非平凡情形只能有
\[
n=2.
\]
又不能有实特征值,否则其特征向量张成一维实不变子空间.
故特征多项式是不可约二次多项式,在 $\mathbb{C}$ 上分裂为两个不同的共轭根.
因此复化后的矩阵有两个互异特征值,故可对角化.
\end{solution}


%例题6.3.5
\begin{example}
设
\[
A=
\begin{pmatrix}
1&-3&-1\\
2&1&0\\
3&1&1
\end{pmatrix}.
\]
证明:
\begin{enumerate}
    \item $A$ 在复数域上可对角化;
    \item $A$ 在有理数域上不可对角化.
\end{enumerate}
\end{example}
\exampleblank

\begin{solution}
特征多项式为
\[
p(\lambda)
=
\lambda^3-3\lambda^2+12\lambda-8.
\]
有
\[
(p,p')=1,
\]
所以在 $\mathbb{C}$ 上三个根互不相同,因此 $A$ 可对角化.

另一方面,由有理根定理,$p$ 的有理根只能是
\[
\pm1,\pm2,\pm4,\pm8,
\]
逐一代入均不为零,所以 $p$ 在 $\mathbb{Q}$ 上没有一次因子.
若 $A$ 能在 $\mathbb{Q}$ 上对角化,其特征值必须都属于 $\mathbb{Q}$,
矛盾.
\end{solution}


%例题6.3.6
\begin{example}
设
\[
A=
\begin{pmatrix}
2&0&0\\
a&2&0\\
b&c&-1
\end{pmatrix}.
\]
\begin{enumerate}
    \item 求所有可能的 Jordan 标准形;
    \item 求 $A$ 可对角化的充分必要条件.
\end{enumerate}
\end{example}
\exampleblank

\begin{solution}
特征值为
\[
2\quad(2\text{ 重}),\qquad -1.
\]
对特征值 $2$,
\[
A-2I=
\begin{pmatrix}
0&0&0\\
a&0&0\\
b&c&-3
\end{pmatrix}.
\]
若 $a=0$,其核维数为 $2$,故 $2$ 对应两个一阶 Jordan 块;
若 $a\neq0$,核维数为 $1$,故对应一个二阶 Jordan 块.
因此只有两种可能:
\[
\boxed{
\diag(2,2,-1)
}
\]
或
\[
\boxed{
J_2(2)\oplus(-1).
}
\]
从而
\[
\boxed{
A\text{ 可对角化}\iff a=0.
}
\]
\end{solution}


%例题6.3.7
\begin{example}
设
\[
f_A(\lambda)
=
\prod\limits_{i=1}^{s}p_i(\lambda),
\]
其中 $p_i$ 是数域 $K$ 上两两不同的首一不可约多项式.
证明:$A$ 的有理标准形只有一个非平凡 Frobenius 块,并且
$A$ 在复数域上可对角化.
\end{example}
\exampleblank

\begin{solution}
特征多项式没有重复不可约因子,因此
\[
m_A=f_A.
\]
于是 $A$ 是循环矩阵意义下的循环变换,有理标准形只有一个友矩阵块
\[
C(f_A).
\]
另一方面在 $\mathbb{C}$ 上,各不可约因子继续分裂为一次因子,
而原特征多项式是平方自由的,所以全部复特征值互不相同,
故 $A$ 在 $\mathbb{C}$ 上可对角化.
\end{solution}


%例题6.3.8
\begin{example}
令
\[
V=M_n(\mathbb{C}),
\qquad
\sigma(X)=AX.
\]
证明:$A$ 可对角化,当且仅当 $\sigma$ 可对角化.
\end{example}
\exampleblank

\begin{solution}
按列分解
\[
M_n(\mathbb{C})
\cong
\mathbb{C}^n\oplus\cdots\oplus\mathbb{C}^n
\]
共 $n$ 份.$\sigma$ 在每一份上都由 $A$ 作用,因此其矩阵相似于
\[
\diag(A,\ldots,A).
\]
故 $\sigma$ 可对角化当且仅当每个 $A$ 块可对角化,即
\[
\boxed{\sigma\text{ 可对角化}\iff A\text{ 可对角化}.}
\]
\end{solution}


%例题6.3.9
\begin{example}
设 $A$ 是可对角化复矩阵,在
\[
V=M_n(\mathbb{C})
\]
上定义
\[
\sigma(X)=AX+XA.
\]
判断 $\sigma$ 是否可对角化.
\end{example}
\exampleblank

\begin{solution}
设
\[
P^{-1}AP=\diag(\lambda_1,\ldots,\lambda_n).
\]
在相应矩阵单位基 $E_{ij}$ 下,
\[
\sigma(E_{ij})
=
(\lambda_i+\lambda_j)E_{ij}.
\]
因此全部 $E_{ij}$ 都是 $\sigma$ 的特征向量,
\[
\boxed{\sigma\text{ 可对角化}.}
\]
\end{solution}


%例题6.3.10
\begin{example}
设 $A$ 为可逆复矩阵,最小多项式次数为 $s$.令
\[
B=(b_{ij})_{s\times s},
\qquad
b_{ij}=\operatorname{tr}A^{i+j},
\quad
1\leq i,j\leq s.
\]
证明
\[
A\text{ 可对角化}
\iff
|B|\neq0.
\]
\end{example}
\exampleblank

\begin{solution}
设 $A$ 的互异特征值为
\[
\lambda_1,\ldots,\lambda_r,
\]
代数重数为 $m_1,\ldots,m_r$.
则
\[
\operatorname{tr}A^k
=
\sum\limits_{\ell=1}^{r}m_\ell\lambda_\ell^k.
\]
故
\[
B=V\,D\,V^{\mathsf T},
\]
其中
\[
V_{i\ell}=\lambda_\ell^i,
\qquad
D=\diag(m_1,\ldots,m_r).
\]
于是
\[
\rank B\leq r.
\]

另一方面,若最小多项式中 $\lambda_\ell$ 的指数为 $d_\ell$,则
\[
s=\sum_{\ell=1}^{r}d_\ell\geq r,
\]
且等号成立当且仅当全部 $d_\ell=1$,也就是 $A$ 可对角化.
若 $A$ 可对角化,则 $s=r$;又 $A$ 可逆,所以各 $\lambda_\ell\neq0$,
$V$ 是非奇异 Vandermonde 型矩阵,因而 $|B|\neq0$.
反之若 $|B|\neq0$,则
\[
s=\rank B\leq r\leq s,
\]
故 $r=s$,从而全部 $d_\ell=1$,$A$ 可对角化.
\end{solution}


%例题6.3.11
\begin{example}
设 $A,B$ 为 $n$ 阶实方阵,满足
\[
AB=BA,\qquad A^n=B^n=I,
\qquad
\operatorname{tr}(AB)=n.
\]
证明
\[
\operatorname{tr}A=\operatorname{tr}B.
\]
\end{example}
\exampleblank

\begin{solution}
因为
\[
x^n-1
\]
在特征 $0$ 中无重根,所以 $A,B$ 都可对角化.
又二者可交换,故可在 $\mathbb{C}$ 上同时对角化:
\[
A\sim\diag(\alpha_1,\ldots,\alpha_n),
\qquad
B\sim\diag(\beta_1,\ldots,\beta_n),
\]
其中
\[
|\alpha_i|=|\beta_i|=1.
\]
题设给出
\[
\displaystyle\sum_{i=1}^{n}\alpha_i\beta_i=n.
\]
$n$ 个模为 $1$ 的复数之和的模达到最大值 $n$,只能全部等于 $1$,故
\[
\alpha_i\beta_i=1.
\]
于是 $B$ 与 $A^{-1}$ 有相同对角表示.
由于 $A$ 为实矩阵,
\[
\operatorname{tr}(A^{-1})
=
\overline{\operatorname{tr}A}
=
\operatorname{tr}A.
\]
所以
\[
\boxed{\operatorname{tr}A=\operatorname{tr}B}.
\]
\end{solution}


%例题6.3.12
\begin{example}
设 $A\in M_n(K)$ 满足
\[
\rank(A-I)+\rank(A+2I)=n.
\]
\begin{enumerate}
    \item 证明
    \[
    A^2+A-2I=O;
    \]
    \item 若 $1$ 是 $A$ 的特征值,$x_0$ 是相应特征向量,证明
    \[
    (I-A)x=x_0
    \]
    无解.
\end{enumerate}
\end{example}
\exampleblank

\begin{solution}
由秩--零化度,
\[
\dim\ker(A-I)+\dim\ker(A+2I)=n.
\]
两个核的交为零,因此
\[
V=\ker(A-I)\oplus\ker(A+2I).
\]
于是 $A$ 在这两个子空间上分别等于 $I$ 与 $-2I$,故
\[
(A-I)(A+2I)=O,
\]
即
\[
\boxed{A^2+A-2I=O}.
\]

若存在 $x$ 使
\[
(I-A)x=x_0,
\]
则
\[
x_0\in\operatorname{Im}(I-A)\cap\ker(I-A).
\]
但由上面的直和分解,$A$ 可对角化,所以
\[
\operatorname{Im}(I-A)\cap\ker(I-A)=\{0\},
\]
与 $x_0\neq0$ 矛盾.
\end{solution}


\newpage

\section{幂零矩阵}

\begin{definition}
若存在正整数 $k$ 使
\[
A^k=O,
\]
则称 $A$ 为幂零矩阵.最小的此类 $k$ 称为幂零指数.
\end{definition}

\begin{proposition}[\markstar]
对复矩阵 $A$,下列命题等价:
\begin{enumerate}
    \item $A$ 幂零;
    \item $A$ 的全部特征值都是 $0$;
    \item
    \[
    \operatorname{tr}A^k=0
    \qquad(\forall k\in\mathbb{Z}_+).
    \]
\end{enumerate}
\end{proposition}


%例题6.4.1
\begin{example}
\begin{enumerate}
    \item 设 $\sigma$ 是 $n$ 维线性空间上的线性变换,满足
    \[
    \sigma^{n-1}\neq0,\qquad
    \sigma^n=0.
    \]
    证明存在一组基使 $\sigma$ 的矩阵为 $J_n(0)$;
    \item 若 $M,N$ 均为 $n$ 阶幂零矩阵,且幂零指数都为 $n$,证明
    \[
    M\sim N.
    \]
\end{enumerate}
\end{example}
\exampleblank

\begin{solution}
取 $v$ 使
\[
\sigma^{n-1}v\neq0.
\]
则
\[
\sigma^{n-1}v,\sigma^{n-2}v,\ldots,\sigma v,v
\]
线性无关,因而构成一组基.$\sigma$ 在此基下正是一个 $n$ 阶幂零
Jordan 块.

第二问中 $M,N$ 都相似于 $J_n(0)$,故
\[
\boxed{M\sim N}.
\]
\end{solution}


%例题6.4.2
\begin{example}
设 $\lambda_0$ 是 $A$ 的 $n$ 重特征值,且
\[
\rank(\lambda_0I-A)=n-1.
\]
\begin{enumerate}
    \item 求使
    \[
    (\lambda_0I-A)^m=0
    \]
    的最小正整数 $m$;
    \item 证明 $(\lambda_0I-A)^{n-1}$ 中必有一个非零列向量是
    $A$ 属于 $\lambda_0$ 的特征向量.
\end{enumerate}
\end{example}
\exampleblank

\begin{solution}
令
\[
N=\lambda_0I-A.
\]
$N$ 为幂零矩阵,且
\[
\rank N=n-1.
\]
这说明 $N$ 的 Jordan 块只有一个,因此
\[
N\sim J_n(0),
\]
其幂零指数为
\[
\boxed{m=n}.
\]
又
\[
N^{n-1}\neq0,\qquad N^n=0.
\]
故 $N^{n-1}$ 至少有一个非零列 $v$,且
\[
Nv=0.
\]
即
\[
(\lambda_0I-A)v=0,
\]
所以 $v$ 是属于 $\lambda_0$ 的特征向量.
\end{solution}


%例题6.4.3
\begin{example}
设 $A,B,C$ 为复矩阵,满足
\[
C=AB-BA,
\qquad
AC=CA.
\]
证明 $C$ 是幂零矩阵.
\end{example}
\exampleblank

\begin{solution}
由 $AC=CA$,$A$ 与 $C^k$ 对任意 $k$ 都可交换.
于是
\[
\begin{aligned}
\operatorname{tr}C^k
&=
\operatorname{tr}\left(C^{k-1}(AB-BA)\right)\\
&=
\operatorname{tr}(C^{k-1}AB)
-\operatorname{tr}(C^{k-1}BA)\\
&=
\operatorname{tr}(BC^{k-1}A)
-\operatorname{tr}(AC^{k-1}B)\\
&=0.
\end{aligned}
\]
故
\[
\operatorname{tr}C^k=0
\qquad(\forall k\geq1).
\]
由幂零判别命题,
\[
\boxed{C\text{ 为幂零矩阵}.}
\]
\end{solution}


%例题6.4.4
\begin{example}
设 $\sigma$ 是 $n$ 维线性空间上的幂零变换,幂零指数为 $p$.
证明:
\begin{enumerate}
    \item
    \[
    p\leq1+\dim\sigma(V);
    \]
    \item 若 $\sigma$ 有两个线性无关的特征向量,则
    \[
    p<n.
    \]
\end{enumerate}
\end{example}
\exampleblank

\begin{solution}
取 $v$ 使
\[
\sigma^{p-1}v\neq0.
\]
则
\[
\sigma v,\sigma^2v,\ldots,\sigma^{p-1}v
\]
线性无关,并且都属于 $\sigma(V)$,故
\[
p-1\leq\dim\sigma(V).
\]

幂零变换的唯一特征值为 $0$,所以特征向量都在
\[
\ker\sigma.
\]
若有两个线性无关特征向量,则
\[
\dim\ker\sigma\geq2,
\]
故
\[
\rank\sigma\leq n-2.
\]
由第一问
\[
p\leq1+\rank\sigma\leq n-1<n.
\]
\end{solution}


%例题6.4.5
\begin{example}
设 $A$ 是 $n$ 阶实矩阵,全部特征值均为实数.
若所有一阶主子式之和为 $0$,所有二阶主子式之和也为 $0$,
证明 $A$ 为幂零矩阵.
\end{example}
\exampleblank

\begin{solution}
设特征值为
\[
\lambda_1,\ldots,\lambda_n\in\mathbb{R}.
\]
一阶主子式之和为特征多项式的第一基本对称系数,故
\[
\displaystyle\sum_i\lambda_i=0.
\]
二阶主子式之和为
\[
\displaystyle\sum_{i<j}\lambda_i\lambda_j=0.
\]
于是
\[
0=
\left(\sum_i\lambda_i\right)^2
=
\sum_i\lambda_i^2
+
2\sum_{i<j}\lambda_i\lambda_j
=
\sum_i\lambda_i^2.
\]
所有 $\lambda_i$ 都为实数,所以
\[
\lambda_i=0
\]
对所有 $i$ 成立.故 $A$ 全部特征值为零,从而幂零.
\end{solution}


%例题6.4.6
\begin{example}
在 $M_n(F)$ 上定义
\[
\sigma_A(B)=AB-BA.
\]
证明:
\begin{enumerate}
    \item $\sigma_A$ 是线性变换;
    \item 若 $A$ 幂零,则 $\sigma_A$ 幂零;
    \item 若 $A$ 为对角矩阵,则 $\sigma_A$ 可对角化.
\end{enumerate}
\end{example}
\exampleblank

\begin{solution}
第一问由矩阵乘法的分配律直接得到.

令
\[
L_A(B)=AB,\qquad R_A(B)=BA.
\]
则
\[
\sigma_A=L_A-R_A,
\qquad
L_AR_A=R_AL_A.
\]
若 $A^m=0$,则
\[
L_A^m=R_A^m=0.
\]
由于二者可交换,二项式展开可知
\[
(L_A-R_A)^{2m-1}=0,
\]
所以 $\sigma_A$ 幂零.

若
\[
A=\diag(a_1,\ldots,a_n),
\]
则矩阵单位 $E_{ij}$ 满足
\[
\sigma_A(E_{ij})
=
(a_i-a_j)E_{ij}.
\]
因此全部 $E_{ij}$ 构成特征向量基,
\[
\boxed{\sigma_A\text{ 可对角化}.}
\]
\end{solution}


%例题6.4.7
\begin{example}
设 $n$ 为奇数,$A,B\in M_n(\mathbb{C})$ 且
\[
A^2=0.
\]
证明
\[
AB-BA
\]
不可逆.
\end{example}
\exampleblank

\begin{solution}
令
\[
C=AB-BA.
\]
直接计算
\[
AC+CA
=
A(AB-BA)+(AB-BA)A
=
A^2B-BA^2
=
0.
\]
因此
\[
AC=-CA.
\]
对任意奇数 $k$,$C^{k-1}$ 与 $A$ 可交换,故
\[
\operatorname{tr}C^k
=
\operatorname{tr}\left(C^{k-1}[A,B]\right)
=
0.
\]
于是 $C$ 的全部奇次幂迹都为零.由 Newton 恒等式,特征多项式只含与
$n$ 同奇偶性的幂次;当 $n$ 为奇数时其常数项必为零.因此
\[
\boxed{\det C=0}.
\]
即 $AB-BA$ 不可逆.
\end{solution}


\newpage

\section{方阵函数}

设 $A=PJP^{-1}$,其中
\[
J=\diag(J_{r_1}(\lambda_1),\ldots,J_{r_s}(\lambda_s)).
\]
对多项式或在谱附近具有足够阶导数的函数 $f$,
\[
f(A)=Pf(J)P^{-1},
\]
而对单个 Jordan 块
\[
J_r(\lambda)=\lambda I+N
\]
有
\[
f(J_r(\lambda))
=
\sum\limits_{j=0}^{r-1}
\dfrac{f^{(j)}(\lambda)}{j!}N^j.
\]

\begin{theorem}[\marktwostar]
若 $A$ 的最小多项式为
\[
m_A(\lambda)
=
\prod_i(\lambda-\lambda_i)^{r_i},
\]
则 $f(A)$ 完全由谱值组
\[
\left\{
f^{(j)}(\lambda_i):
0\leq j\leq r_i-1
\right\}
\]
决定.反过来,对任意给定的这样一组数,都存在多项式
$f_*(\lambda)$ 具有这些 Hermite 插值数据,并满足
\[
f(A)=f_*(A).
\]
\end{theorem}

\begin{proposition}[\markstar]
当谱位于相应收敛圆内时,可以用幂级数定义
\[
e^A,\qquad
\sin A,\qquad
\cos A,\qquad
\log(I+A),\qquad
(I+A)^a.
\]
特别地,
\[
e^A=\sum\limits_{k=0}^{\infty}\dfrac{A^k}{k!}.
\]
\end{proposition}


%例题6.5.1
\begin{example}
设
\[
A=
\begin{pmatrix}
0&1&1&1\\
0&0&1&1\\
0&0&0&1\\
0&0&0&0
\end{pmatrix}.
\]
求 $A$ 的 Jordan 标准形,并计算 $e^A$.
\end{example}
\exampleblank

\begin{solution}
有
\[
A^3\neq0,\qquad A^4=0,
\]
故幂零指数为 $4$,于是
\[
\boxed{A\sim J_4(0)}.
\]
由于 $A^4=0$,
\[
e^A
=
I+A+\dfrac{A^2}{2}+\dfrac{A^3}{6}.
\]
计算得
\[
\boxed{
e^A=
\begin{pmatrix}
1&1&\dfrac32&\dfrac{13}{6}\\
0&1&1&\dfrac32\\
0&0&1&1\\
0&0&0&1
\end{pmatrix}.
}
\]
\end{solution}


%例题6.5.2
\begin{example}
计算下列方阵函数:
\begin{enumerate}
    \item
    \[
    \sqrt A,\qquad
    A=
    \begin{pmatrix}
    3&1\\
    -1&5
    \end{pmatrix};
    \]
    \item
    \[
    e^A,\qquad
    A=
    \begin{pmatrix}
    4&-2\\
    6&-3
    \end{pmatrix};
    \]
    \item
    \[
    \sin A,\qquad
    A=
    \begin{pmatrix}
    \pi-1&1\\
    -1&\pi+1
    \end{pmatrix};
    \]
    \item
    \[
    \sin A,\qquad
    A=
    \begin{pmatrix}
    3&-1&-1\\
    1&1&-1\\
    0&0&2
    \end{pmatrix}.
    \]
\end{enumerate}
\end{example}
\exampleblank

\begin{solution}
第一问写
\[
A=4I+N,
\qquad
N^2=0.
\]
对 $f(x)=\sqrt x$,
\[
f(4)=2,\qquad f'(4)=\dfrac14.
\]
故
\[
\boxed{
\sqrt A
=
2I+\dfrac14N
=
\begin{pmatrix}
\dfrac74&\dfrac14\\[2mm]
-\dfrac14&\dfrac94
\end{pmatrix}.
}
\]

第二问 $A$ 的特征值为 $0,1$,且可对角化,故用一次插值多项式
\[
e^x\equiv1+(e-1)x
\]
作用于 $A$:
\[
\boxed{
e^A=I+(e-1)A.
}
\]

第三问写
\[
A=\pi I+N,
\qquad N^2=0.
\]
于是
\[
\sin A
=
\sin\pi\,I+\cos\pi\,N
=
-N,
\]
即
\[
\boxed{
\sin A=
\begin{pmatrix}
1&-1\\
1&-1
\end{pmatrix}.
}
\]

第四问写
\[
A=2I+N,
\qquad N^2=0.
\]
因此
\[
\boxed{
\sin A
=
\sin2\,I+\cos2\,(A-2I).
}
\]
\end{solution}


%例题6.5.3
\begin{example}
设
\[
A=
\begin{pmatrix}
0&1&0\\
0&0&0\\
0&0&0
\end{pmatrix},
\qquad
B=
\begin{pmatrix}
0&0&1\\
0&0&0\\
0&0&0
\end{pmatrix}.
\]
问是否存在三阶复方阵 $X$ 以及多项式 $f,g\in\mathbb{C}[x]$,使
\[
A=f(X),\qquad B=g(X).
\]
\end{example}
\exampleblank

\begin{solution}
若存在,则 $X$ 必同时与 $A,B$ 可交换.
直接解
\[
XA=AX,\qquad XB=BX
\]
可得所有公共可交换矩阵恰为
\[
X=
aI+bA+cB.
\]
并且
\[
(X-aI)^2=0.
\]
因此任意多项式 $h(X)$ 都只能写成
\[
h(X)=uI+v(bA+cB),
\]
也就是说其非数量部分最多是一维的.

但 $A,B$ 线性无关,不可能同时属于同一个一维非数量子空间.
故
\[
\boxed{\text{不存在这样的 }X,f,g.}
\]
\end{solution}


%例题6.5.4
\begin{example}
证明对任意复方阵 $A$,
\[
e^A
\]
都可逆.
\end{example}
\exampleblank

\begin{solution}
因为 $A$ 与 $-A$ 可交换,幂级数乘法给出
\[
e^Ae^{-A}=e^{A-A}=I.
\]
故
\[
\boxed{(e^A)^{-1}=e^{-A}}.
\]
\end{solution}


%例题6.5.5
\begin{example}
求
\[
|e^A|,
\]
其中 $A$ 为 $n$ 阶复方阵.
\end{example}
\exampleblank

\begin{solution}
若 $A$ 的特征值为
\[
\lambda_1,\ldots,\lambda_n
\]
(重数计入),则 $e^A$ 的特征值为
\[
e^{\lambda_1},\ldots,e^{\lambda_n}.
\]
因此
\[
\boxed{
\det(e^A)
=
e^{\lambda_1+\cdots+\lambda_n}
=
e^{\operatorname{tr}A}.
}
\]
\end{solution}


\newpage

\section{矩阵可交换问题}

设
\[
A=PJP^{-1},
\qquad
J=\diag(J_1,\ldots,J_s),
\qquad
J_i=J_{k_i}(\lambda_i).
\]
令
\[
B=PB_1P^{-1},
\qquad
B_1=(B_{ij}).
\]
则
\[
AB=BA
\iff
J_iB_{ij}=B_{ij}J_j.
\]
当 $\lambda_i\neq\lambda_j$ 时必有 $B_{ij}=0$;
当 $\lambda_i=\lambda_j$ 时,$B_{ij}$ 是相应的截断 Toeplitz 型矩阵.

\begin{theorem}[\marktwostar]
与复矩阵 $A$ 可交换的全部矩阵,可通过上述 Jordan 分块逐块描述.
特别地,当
\[
m_A=f_A,
\]
即 $A$ 为循环变换时,其交换子代数恰为
\[
\boxed{\mathbb{C}[A]}.
\]
\end{theorem}


%例题6.6.1
\begin{example}
求与下列矩阵可交换的全部方阵:
\begin{enumerate}
    \item
    \[
    A_1=
    \begin{pmatrix}
    3&2&-5\\
    2&6&-10\\
    1&2&-3
    \end{pmatrix};
    \]
    \item
    \[
    A_2=
    \begin{pmatrix}
    6&6&-15\\
    1&5&-5\\
    1&2&-2
    \end{pmatrix};
    \]
    \item
    \[
    A_3=
    \begin{pmatrix}
    14&-2&-7&-1\\
    20&-2&-11&-2\\
    19&-3&-9&-1\\
    -6&1&3&1
    \end{pmatrix}.
    \]
\end{enumerate}
\end{example}
\exampleblank

\begin{solution}
第一问可取
\[
P_1=
\begin{pmatrix}
1&1&-2\\
2&0&1\\
1&0&0
\end{pmatrix},
\qquad
P_1^{-1}A_1P_1
=
J_2(2)\oplus(2).
\]
因此全部交换矩阵为
\[
\boxed{
B=P_1
\begin{pmatrix}
a&b&c\\
0&a&0\\
0&d&e
\end{pmatrix}
P_1^{-1},
\qquad
a,b,c,d,e\in\mathbb{C}.
}
\]

第二问可取
\[
P_2=
\begin{pmatrix}
3&1&-2\\
1&0&1\\
1&0&0
\end{pmatrix},
\qquad
P_2^{-1}A_2P_2
=
J_2(3)\oplus(3).
\]
故
\[
\boxed{
B=P_2
\begin{pmatrix}
a&b&c\\
0&a&0\\
0&d&e
\end{pmatrix}
P_2^{-1}.
}
\]

第三问可取
\[
P_3=
\begin{pmatrix}
2&13&1&1\\
3&20&0&3\\
3&19&0&1\\
-1&-6&0&0
\end{pmatrix},
\]
且
\[
P_3^{-1}A_3P_3
=
J_3(1)\oplus(1).
\]
因此
\[
\boxed{
B=P_3
\begin{pmatrix}
a&b&c&d\\
0&a&b&0\\
0&0&a&0\\
0&0&e&f
\end{pmatrix}
P_3^{-1},
}
\]
其中参数任意.
\end{solution}


%例题6.6.2
\begin{example}
证明:若方阵 $A,B$ 可交换,则 $A$ 的每个根子空间都是 $B$ 的不变子空间.
\end{example}
\exampleblank

\begin{solution}
设 $\lambda$ 是 $A$ 的特征值,根子空间为
\[
W_\lambda=\ker(A-\lambda I)^m
\]
其中 $m$ 充分大.由 $AB=BA$,
\[
B(A-\lambda I)^m
=
(A-\lambda I)^mB.
\]
若 $x\in W_\lambda$,则
\[
(A-\lambda I)^mBx
=
B(A-\lambda I)^mx
=
0.
\]
故
\[
Bx\in W_\lambda.
\]
因此
\[
\boxed{W_\lambda\text{ 是 }B\text{ 的不变子空间}.}
\]
\end{solution}
